\section{Breaking Objects}
  There are two main ways of breaking objects.
  You can deal damage to objects with attacks, similarly to how you can deal damage to creatures.
  Alternately, you can attempt to sunder the object with the Might skill (see \pcref{Might}).

  \subsection{Damaging Objects}
    Objects have \glossterm{hit points} like creatures.
    They can also have \glossterm{hardness}, which reduces damage they take.
    Whenever an object takes damage, it first subtracts its hardness from that damage before applying that damage to its hit points.
    Objects are not normally subject to \glossterm{critical hits}.
    When objects lose enough hit points, they can become \glossterm{broken} (see \pcref{Broken and Destroyed Objects}).

    \subsubsection{Vulnerable and Resistant Objects}
      An object that is \resistant to an attack doubles its \glossterm{hardness} against that attack.
      Likewise, an object that is \vulnerable to an attack has no hardness against that attack.

  \subsection{Object Statistics}
    An object's size primarily influences the number of \glossterm{hit points} it has.
    The primary material it is constructed from determines its \glossterm{hardness}, and can modify the number of hit points it has.
    Details are given in \tref{Object Statistics By Size} and \tref{Object Statistics By Material}.

    Oddly shaped objects, such as a ladder or pillar, should generally use some size category in between their largest and smallest dimensions.
    Consider these examples:
    \begin{raggeditemize}
      \item A 20 foot cube would have the hit points of a Huge object.
      \item A stone pillar that is 20 feet wide, 20 feet long, and 100 feet tall would have the hit points of a Huge object.
      \item A 20 foot tall, 20 foot wide wall of typical thickness would have the hit points of a Large object.
    \end{raggeditemize}

    These rules are more detailed than you should really need.
    During a typical game session, it's often best to just guess whether a character could plausibly sunder or smash an object rather than consulting these tables.

    % TODO: Should object HP be determined by size or by weight?
    \begin{dtable}
      \lcaption{Object Statistics By Size}
      \begin{dtabularx}{\textwidth}{l X X}
        \tb{Size}  & \tb{Hit Points} & \tb{Sunder DV Modifier} \tableheaderrule
        Fine       & 2               & 0  \\
        Diminutive & 5               & 0  \\
        Tiny       & 10              & 0  \\
        Small      & 25              & \plus5  \\
        Medium     & 50              & \plus10 \\
        Large      & 100             & \plus15 \\
        Huge       & 250             & \plus20 \\
        Gargantuan & 500             & \plus25 \\
        Colossal   & 1000            & \plus30 \\
      \end{dtabularx}
    \end{dtable}

    \begin{dtable}
      \lcaption{Object Statistics By Material}
      \begin{dtabularx}{\textwidth}{l X X X}
        \tb{Material}   & \tb{Hardness}\fn{1} & \tb{HP Multiplier}\fn{2} \tableheaderrule
        Adamantine      & 30            & \mult3   \\
        Glass           & 5             & \mult1/2 \\
        Ice             & 0             & \mult1/2 \\
        Iron or steel   & 15            & \mult2   \\
        Leather or hide & 5             & \tdash   \\
        Mithral         & 20            & \mult2   \\
        Packed earth    & 5             & \tdash   \\
        Paper or cloth  & 0             & \mult1/2 \\
        Rope            & 5             & \tdash   \\
        Stone           & 10            & \mult2   \\
        Wood            & 5             & \tdash   \\
      \end{dtabularx}
      1. An object's \glossterm{hardness} also increases the difficulty value of checks to sunder it with raw Strength. \\
      2. Any value here modifies the number of hit points the object would normally have based on its size.
    \end{dtable}

  \subsection{Broken and Destroyed Objects}\label{Broken and Destroyed Objects}
    If an object would gain its first \glossterm{vital wound}, typically by reaching negative hit points, it becomes \glossterm{broken} instead.
    Unlike creatures, objects do not reset their negative hit points to 0 at the start of their turn.
    Instead, they continue to track their negative hit points indefinitely.
    If an object's negative hit points reach ten times its maximum hit points, it becomes \glossterm{destroyed}.

    \parhead{Broken Objects}\nonsectionlabel{Broken Objects}
    Broken objects cannot be used for their intended purpose, but still retain enough of their original form to be repaired without too much work.
    For example, a broken wall lies in pieces on the ground and no longer blocks passage, but can be repaired with far less effort than would be required to create a wall from scratch.
    Magic items that are broken retain their magical properties once fixed.
    Broken (but not destroyed) objects can be repaired with the Craft skill (see \pcref{Craft}).

    Most magical effects that create temporary objects are fully destroyed when they become broken.

    \parhead{Destroyed Objects}\nonsectionlabel{Destroyed Objects}
    Destroyed objects have been damaged beyond hope of any sort of repair short of crafting the object again from raw materials.
    For example, a destroyed wall is reduced to dust or small, useless chunks of rubble.
    Magic items that are destroyed irrevocably lose their magical properties.
    The remains of a destroyed object generally occupy a space one size category smaller than the original object.
    Destroyed objects can be rebuilt with the Craft skill, but it requires significant time and investment.

  \subsection{Breaking Equipment}\label{Breaking Equipment}
    Normally, a character's equipment cannot be damaged or otherwise affected by attacks.
    This includes worn items, anything held in your hands, and anything in a secure storage like a small backpack.
    Such items are considered \glossterm{attended}.
    They are unaffected by damage caused by area effects, and cannot be targeted individually.
    Some abilities can specifically target \glossterm{attended} objects, as indicated in their descriptions.

    \subsubsection{Loose Equipment}\label{Loose Equipment}
      Some items are explicitly \glossterm{loose equipment}.
      Loose equipment does not gain the protections listed above while worn as equipment.
      It can be individually targeted by attackers, and is affected by area effects just like any other object in the area.
